\begin{abstract}
Are today's vehicle-to-everything (V2X) and fifth-generation (5G) communication
paths ready for tomorrow's connected and automated vehicle (CAV) applications?
Conventional connected-vehicle (CV) and intelligent transportation system (ITS)
applications exchange compact, stateless SAE International J2735 messages.
Complex CAV applications instead use correlated, stateful operations. These
operations may carry larger objects and generate event-triggered bursts or
sustained streams in either direction. The current landscape lacks a common
interface for both classes. Existing studies do not establish where either path
misses application deadlines when carrying both workload classes.

In this paper, we contribute the Intersection Programming Interface (IPI) and a
detailed real-vehicle communication-latency study. IPI provides a
transport-independent interface for the SAE International J2735 message set and
correlated CAV operations. The study evaluates both classes over a commercial
Long-Term Evolution cellular V2X (LTE C-V2X) direct PC5 path and a private 5G
New Radio (NR) network-assisted Uu path. The 5G deployment uses one physical CAV
user equipment (UE), a dedicated radio, and a clean 40-MHz n48 channel without
ambient contention. Competing traffic and application-client demand are
introduced only as experimental variables. Experiments vary payload size,
direction, transport, radio conditions, and interruption. We report
response round-trip time (RTT), the percentage of issued requests answered by
each deadline, and directional application goodput.

The measurements yield three insights. 1) The measured commercial direct V2X
path supports compact J2735 messages only within a limited payload and coverage
envelope. 2) The CAV application-packet size that the measured 5G uplink
completes within a decision deadline is severely limited. 3) 5G and
sixth-generation (6G) radios must support event-triggered bursts and sustained
streams in either direction while isolating concurrent flows if cellular
networks are to support CAV and ITS applications that advance the goal of zero
road fatalities.
\end{abstract}
